%============================================================
% CH1 · P1 · PROBLEM 1 — THIRD MAXIMUM NUMBER (LeetCode 414)
%============================================================
\subsection{Problem 1.1.1: Third Maximum Number}
\problemlabel{1.1.1}{ssec:ch1-third-maximum}

\problemcard{1.1.1}{Third Maximum Number}{LeetCode 414}{https://leetcode.com/problems/third-maximum-number/}{\exeasy\ \lc{Easy}}{Single pass keeping the three largest distinct values}

\subsubsection*{1. Problem statement}

You are given an integer array \texttt{nums}. Return the \textbf{third largest distinct value} in it. If the array contains fewer than three distinct values, return the \textbf{largest} value instead. ``Distinct'' matters: in $[2, 2, 3, 1]$ the two 2's count as one value, so the distinct values are $3 > 2 > 1$ and the answer is $1$.

\textbf{Constraints:} $1 \le n \le 10^4$, \; $-2^{31} \le \texttt{nums}[i] \le 2^{31} - 1$ (the whole \texttt{int} range). \textbf{Follow-up:} can you do it in $\bigO{n}$ time?

\begin{table}[H]
\centering\small
\begin{tabular}{@{}llcl@{}}
\toprule
& \textbf{nums} & \textbf{Output} & \textbf{Why} \\
\midrule
Example 1 & $[3, 2, 1]$ & 1 & distinct values $3 > 2 > 1$: the third is $1$ \\
Example 2 & $[1, 2]$ & 2 & only two distinct values: return the maximum \\
Example 3 & $[2, 2, 3, 1]$ & 1 & distinct values $3 > 2 > 1$; the second 2 does not count \\
\bottomrule
\end{tabular}
\end{table}

\begin{diagrambox}
\begin{tikzpicture}
    \begin{scope}
        \dparray{e}{0}{Ex.\ 1}{3/acell, 2/acell, 1/anscell}
    \end{scope}
    \begin{scope}[xshift=4.3cm]
        \dparray{f}{0}{Ex.\ 2}{1/acell, 2/anscell}
    \end{scope}
    \begin{scope}[xshift=7.6cm]
        \dparray{g}{0}{Ex.\ 3}{2/acell, 2/aout, 3/acell, 1/anscell}
    \end{scope}
    \node[note, text width=12cm, anchor=north] at (5.6,-0.55) {orange: the answer; grey: a duplicate that must be ignored. In Example 2 there is no third value, so the answer falls back to the maximum.};
\end{tikzpicture}
\figcap{The three official examples.}
\end{diagrambox}

\subsubsection*{2. The problem model}
\begin{modelbox}[Problem Model --- Third Maximum Number]
\renewcommand{\arraystretch}{1.35}
\begin{tabularx}{\linewidth}{@{}>{\bfseries}l X@{}}
Input & one array of $n$ integers, any value of the \texttt{int} range (duplicates allowed) \\
Computed & one \emph{value}: the third element of the \emph{set} of values sorted decreasingly, or its first element if the set has fewer than three elements \\
Over & the distinct values (the positions do not matter at all) \\
Sizes & $n \le 10^4$: even $\bigO{n^2} = 10^8$ would pass; the follow-up asks for $\bigO{n}$ \\
Trap & values reach $-2^{31}$, so no \texttt{int} value is free to mean ``slot still empty'' \\
\end{tabularx}
\end{modelbox}

\begin{diagrambox}
\begin{tikzpicture}
    \dparray{m}{0}{nums}{4/acell, 9/acell, 4/acell, 7/acell, 1/acell, 9/acell, 6/acell}
    \draw[flow] (6.5,0) -- node[lbl, above=2pt]{set} (8.1,0);
    \begin{scope}[xshift=8.3cm]
        \node[cell, fill=basefill] at (0.45,0) {9};
        \node[cell, fill=basefill] at (1.35,0) {7};
        \node[anscell] at (2.25,0) {6};
        \node[cell, fill=lightgray, text=midgray] at (3.15,0) {4};
        \node[cell, fill=lightgray, text=midgray] at (4.05,0) {1};
        \node[font=\scriptsize, text=darkgray] at (0.45,-0.6) {1st}; \node[font=\scriptsize, text=darkgray] at (1.35,-0.6) {2nd};
        \node[font=\scriptsize\bfseries, text=accentorange] at (2.25,-0.6) {3rd};
    \end{scope}
\end{tikzpicture}
\figcap{The model: forget positions and multiplicities, keep the set of values. Only its three largest members can ever matter. (This array is the trace example of Section~5.)}
\end{diagrambox}

\subsubsection*{3. Key observations}

\begin{itemize}[leftmargin=*]
    \item \textbf{Only the top three distinct values matter.} A value smaller than three other distinct values can never become the answer, no matter what follows. So a state of three numbers is enough --- Pattern~1 with $k = 3$ distinct.
    \item \textbf{Sorting without removing duplicates is wrong.} The tempting ``sort decreasingly, return \texttt{nums[2]}'' gives $2$ on Example 3: the sorted array $3, 2, 2, 1$ has a 2 at index 2. Duplicates must be skipped.
    \item \textbf{A duplicate must be skipped before it is compared.} Without the skip, the second 9 in $[4, 9, 4, 7, 1, 9, 6]$ would be ``smaller than $m_1 = 9$, larger than $m_2$'' and push a 9 into the second slot.
    \item \textbf{\texttt{INT\_MIN} is a legal value.} If ``empty slot'' is encoded as \texttt{INT\_MIN}, the input $[1, 2, -2^{31}]$ ends with $m_3 = -2^{31}$, which looks empty, and the code returns $2$ instead of $-2147483648$.
\end{itemize}

\begin{diagrambox}
\begin{tikzpicture}
    \begin{scope}
        \dparray{w}{0}{sorted}{3/acell, 2/acell, 2/aread, 1/acell}
        \aptr{w2}{index 2}{warningred}
        \node[note, anchor=west, text width=3cm, align=left] at (3.8,0) {returns $2$ --- \wa{}\\(the answer is $1$)};
    \end{scope}
    \begin{scope}[xshift=7.6cm]
        \dparray{v}{0}{slots}{2/acell, 1/acell, {\tiny INT\_MIN}/aread}
        \node[note, anchor=north, text width=3.6cm] at (1.35,-0.55) {input $[1, 2, -2^{31}]$: is $m_3$ empty or real? Code returns $2$ --- \wa{}};
    \end{scope}
\end{tikzpicture}
\figcap{Two tempting wrong solutions. Left: sorting without deduplication on Example 3. Right: \texttt{int} slots with \texttt{INT\_MIN} meaning ``empty''.}
\end{diagrambox}

\subsubsection*{4. Approaches}

\begin{bruteforcebox}
\textbf{Idea.} Sort decreasingly, remove adjacent duplicates with \texttt{unique}, and read index 2 (or index 0 if fewer than three values remain).
\begin{lstlisting}[style=cpp]
#include <bits/stdc++.h>
using namespace std;

class Solution {
public:
    int thirdMax(vector<int>& nums) {
        vector<int> v = nums;                                  // keep the input intact
        sort(v.begin(), v.end(), greater<int>());              // largest first
        v.erase(unique(v.begin(), v.end()), v.end());          // drop adjacent duplicates
        return v.size() >= 3 ? v[2] : v[0];
    }
};

int main() {
    vector<vector<int>> tests = {{3, 2, 1}, {1, 2}, {2, 2, 3, 1}};
    for (auto& t : tests) cout << Solution().thirdMax(t) << " ";   // 1 2 1
    cout << "\n";
}
\end{lstlisting}
\begin{complexitybox}
\textbf{Time:} $\bigO{n \log n}$ for the sort --- about $1.4\cdot10^5$ comparisons for $n = 10^4$, so \ac{}. \textbf{Space:} $\bigO{n}$ for the copy. It is correct and short, but it sorts $n$ values to learn three of them, and it does not meet the $\bigO{n}$ follow-up.
\end{complexitybox}
\end{bruteforcebox}

\begin{algorithmbox}
\textbf{Idea.} Scan once and keep the three largest distinct values seen so far, $a > b > c$, in \texttt{long long} variables initialised to \texttt{LLONG\_MIN} (a value no \texttt{int} can take, so it safely means ``empty''). For each $x$:
\begin{enumerate}[leftmargin=*]
    \item if $x$ equals $a$, $b$ or $c$, skip it (already counted);
    \item if $x > a$: shift down ($c \gets b$, $b \gets a$) and set $a \gets x$;
    \item else if $x > b$: $c \gets b$, $b \gets x$;
    \item else if $x > c$: $c \gets x$.
\end{enumerate}
At the end, $c$ is still empty exactly when fewer than three distinct values exist; then return $a$.
\begin{lstlisting}[style=cpp]
#include <bits/stdc++.h>
using namespace std;

class Solution {
public:
    int thirdMax(vector<int>& nums) {
        long long a = LLONG_MIN, b = LLONG_MIN, c = LLONG_MIN;   // a > b > c, LLONG_MIN = empty
        for (int x : nums) {
            if (x == a || x == b || x == c) continue;           // duplicate value: skip
            if (x > a)      { c = b; b = a; a = x; }            // new largest: shift down twice
            else if (x > b) { c = b; b = x; }                   // new second: shift once
            else if (x > c) { c = x; }                          // new third
        }
        return c == LLONG_MIN ? (int)a : (int)c;                // < 3 distinct: the maximum
    }
};

int main() {
    vector<vector<int>> tests = {{3, 2, 1}, {1, 2}, {2, 2, 3, 1}, {1, 2, INT_MIN}};
    for (auto& t : tests) cout << Solution().thirdMax(t) << " ";   // 1 2 1 -2147483648
    cout << "\n";
}
\end{lstlisting}
\begin{complexitybox}
\textbf{Time:} $\bigO{n}$ --- at most six comparisons per element. \textbf{Space:} $\bigO{1}$. \textbf{Verdict:} \ac{}, and it answers the follow-up.
\end{complexitybox}
\end{algorithmbox}

\begin{invariantbox}[Why the three slots are correct]
\textbf{Invariant.} Before processing \texttt{nums[i]}, let $D_i$ be the set of distinct values in \texttt{nums[0..i)}. Then $a$, $b$, $c$ are the largest, second largest and third largest elements of $D_i$, where a missing element is represented by \texttt{LLONG\_MIN}.

\textbf{Initialisation.} $D_0 = \emptyset$ and all three slots are empty.

\textbf{Maintenance.} Let $x = \texttt{nums}[i]$. If $x \in \{a, b, c\}$ then $x \in D_i$, so $D_{i+1} = D_i$ and nothing changes --- correct. Otherwise $x$ is either new, or an old value smaller than $c$ (only possible when $c$ is a real value, i.e.\ $|D_i| \ge 3$). The top three of $D_i \cup \{x\}$ are the top three of $\{a, b, c, x\}$, because every other element of $D_i$ is below $c$. The four cases compare $x$ with $a$, $b$, $c$ and insert it at its rank, dropping the old $c$ when $x$ enters the top three; an old value below $c$ falls into ``no change''. An empty slot (\texttt{LLONG\_MIN}) is smaller than every \texttt{int}, so a new value always fills the first empty slot at or after its rank.

\textbf{Termination.} After the last element, $c$ is a real value iff $|D_n| \ge 3$, and then it is the third largest distinct value; otherwise the answer is the maximum $a$, which is real because $n \ge 1$. \hfill$\blacksquare$
\end{invariantbox}

\begin{alternativebox}
\textbf{A set of size at most 3.} \texttt{std::set} removes duplicates and keeps its elements sorted. Insert every value; whenever the set grows to 4 elements, erase the smallest. At the end the set holds the (up to) three largest distinct values.
\begin{lstlisting}[style=cpp]
#include <bits/stdc++.h>
using namespace std;

class Solution {
public:
    int thirdMax(vector<int>& nums) {
        set<int> top;                                  // sorted, no duplicates, size <= 3
        for (int x : nums) {
            top.insert(x);                             // a duplicate is ignored by the set
            if (top.size() > 3) top.erase(top.begin());   // drop the smallest of four
        }
        return top.size() == 3 ? *top.begin() : *top.rbegin();
    }
};

int main() {
    vector<vector<int>> tests = {{3, 2, 1}, {1, 2}, {2, 2, 3, 1}, {1, 2, INT_MIN}};
    for (auto& t : tests) cout << Solution().thirdMax(t) << " ";   // 1 2 1 -2147483648
    cout << "\n";
}
\end{lstlisting}
Each operation is on a set of at most 4 elements, so the total time is $\bigO{n}$ (with a larger constant than the three variables) and the memory is $\bigO{1}$. There is no sentinel at all: emptiness is \texttt{top.size()}. The same code with $3$ replaced by $k$ returns the $k$-th largest distinct value in $\bigO{n \log k}$.

\begin{diagrambox}
\begin{tikzpicture}[x=1cm, y=1cm, sc/.style={cell, minimum width=7mm, minimum height=6mm, font=\small}]
    \foreach \x/\st/\drop [count=\k from 0] in {4/{4}/, 9/{4,9}/, 4/{4,9}/, 7/{4,7,9}/, 1/{4,7,9}/1, 9/{4,7,9}/, 6/{6,7,9}/4}{
        \node[font=\scriptsize\bfseries, text=deepblue] at (\k*1.75,0.6) {$+\x$};
        \foreach \v [count=\j from 0] in \st { \node[sc, fill=basefill] at (\k*1.75,-\j*0.62) {\v}; }
        \ifx\drop\empty\else \node[font=\scriptsize, text=warningred] at (\k*1.75,-2.05) {drop \drop}; \fi
    }
\end{tikzpicture}
\figcap{The set after each insertion on $[4, 9, 4, 7, 1, 9, 6]$ (smallest on top). Inserting $1$ makes four elements and $1$ is dropped at once; inserting $6$ drops $4$. The final smallest element, $6$, is the answer.}
\end{diagrambox}
\end{alternativebox}

\begin{tricksbox}
\begin{itemize}[leftmargin=*]
    \item \textbf{Three passes, no sentinel trouble.} Find the maximum $m_1$; then the maximum among values $< m_1$; then among values $< m_2$. Each pass is a plain maximum with a bound, $3n$ steps in total. If a pass finds nothing, fewer than three distinct values exist. This is the easiest version to prove and generalises to ``$k$ passes for the $k$-th distinct maximum'' ($\bigO{nk}$).
    \item \textbf{A counter instead of a sentinel.} Keep \texttt{int filled} $\in \{0,1,2,3\}$ and compare only with filled slots; then the slots may be \texttt{int}.
    \item \textbf{Why \texttt{long long} works:} \texttt{int} values are promoted when compared with \texttt{long long}, and \texttt{LLONG\_MIN} $< -2^{31}$, so no input can collide with it.
\end{itemize}
\end{tricksbox}

\subsubsection*{5. Trace}

\begin{tracebox}[{One pass on $\texttt{nums} = [4, 9, 4, 7, 1, 9, 6]$}]
\centering
\begin{tikzpicture}
    \begin{scope}
        \dparray{t}{0}{step 2}{4/adone, 9/aread, 4/acell, 7/acell, 1/acell, 9/acell, 6/acell}
        \aptr{t1}{x = 9}{accentorange}
        \node[note, anchor=west, text width=4.2cm, align=left] at (6.6,0) {$9 > a = 4$: shift, $(a,b,c) = (9, 4, -)$};
    \end{scope}
    \begin{scope}[yshift=-1.95cm]
        \dparray{u}{0}{step 3}{4/adone, 9/adone, 4/aout, 7/acell, 1/acell, 9/acell, 6/acell}
        \aptr{u2}{x = 4}{midgray}
        \node[note, anchor=west, text width=4.2cm, align=left] at (6.6,0) {$4 = b$: duplicate, skipped};
    \end{scope}
    \begin{scope}[yshift=-3.9cm]
        \dparray{v}{0}{step 4}{4/adone, 9/adone, 4/aout, 7/aread, 1/acell, 9/acell, 6/acell}
        \aptr{v3}{x = 7}{accentorange}
        \node[note, anchor=west, text width=4.2cm, align=left] at (6.6,0) {$b < 7 < a$: $(a,b,c) = (9, 7, 4)$ --- all three slots full};
    \end{scope}
    \begin{scope}[yshift=-5.85cm]
        \dparray{w}{0}{step 7}{4/adone, 9/adone, 4/aout, 7/adone, 1/aout, 9/aout, 6/anscell}
        \aptr{w6}{x = 6}{accentorange}
        \node[note, anchor=west, text width=4.2cm, align=left] at (6.6,0) {$c < 6 < b$: $c \gets 6$; answer $c = 6$};
    \end{scope}
\end{tikzpicture}
\figcap{Grey cells were skipped ($4$ and $9$ are duplicates, $1$ is below $c$). The answer is the orange cell.}

\medskip
\small
\begin{tabular}{@{}clllcl@{}}
\toprule
\textbf{step} & $x$ & \textbf{case} & $a$ & $b$ & $c$ \\
\midrule
(start) & & & --- & --- & --- \\
1 & 4 & $x > a$ (empty) & 4 & --- & --- \\
2 & 9 & $x > a$ & 9 & 4 & --- \\
3 & 4 & duplicate of $b$ & 9 & 4 & --- \\
4 & 7 & $b < x < a$ & 9 & 7 & 4 \\
5 & 1 & $x < c$ & 9 & 7 & 4 \\
6 & 9 & duplicate of $a$ & 9 & 7 & 4 \\
7 & 6 & $c < x < b$ & 9 & 7 & \textbf{6} \\
\bottomrule
\end{tabular}
\end{tracebox}

\begin{tracebox}[{Example 3: $\texttt{nums} = [2, 2, 3, 1]$}]
\centering
\small
\begin{tabular}{@{}clllcl@{}}
\toprule
\textbf{step} & $x$ & \textbf{case} & $a$ & $b$ & $c$ \\
\midrule
1 & 2 & $x > a$ (empty) & 2 & --- & --- \\
2 & 2 & duplicate of $a$ & 2 & --- & --- \\
3 & 3 & $x > a$ & 3 & 2 & --- \\
4 & 1 & $x > c$ (empty) & 3 & 2 & \textbf{1} \\
\bottomrule
\end{tabular}
\quad
\begin{minipage}{0.4\linewidth}\small
$c = 1$ is real, so the answer is $1$. For Example 2, $[1, 2]$, the states are $(1,-,-)$ and $(2,1,-)$; $c$ stays empty and the answer is $a = 2$.
\end{minipage}
\end{tracebox}

\subsubsection*{6. Test suite}

\begin{testbox}
\centering\small
\begin{tabular}{@{}llc@{}}
\toprule
\textbf{nums} & \textbf{Why it is interesting} & \textbf{Expected} \\
\midrule
$[1]$ & a single element & 1 \\
$[5, 5, 5]$ & one distinct value, $n = 3$ & 5 \\
$[1, 2]$ & official Example 2: fallback to the maximum & 2 \\
$[2, 2, 3, 1]$ & official Example 3: duplicate in the middle & 1 \\
$[3, 3, 2, 2, 1, 1]$ & every value twice & 1 \\
$[1, 2, 2, 5, 3, 5]$ & duplicates of the first and second maximum & 2 \\
$[1, 2, -2^{31}]$ & the answer is \texttt{INT\_MIN} (sentinel trap) & $-2147483648$ \\
$[-2^{31}, -2^{31}, -2^{31}]$ & only \texttt{INT\_MIN}, fallback & $-2147483648$ \\
$[2^{31}-1, 1, -2^{31}]$ & both ends of the \texttt{int} range & $-2147483648$ \\
$[1, 2, \dots, 10^4]$ & largest $n$, increasing: every element is a new maximum & 9998 \\
\bottomrule
\end{tabular}
\end{testbox}
All expected values were produced by the solution and by the brute force; in addition, the one-pass solution, the set version and the three-pass version were compared with the brute force on 900 random arrays ($n \le 9$, values in $[-3, 3]$ or drawn from $\{-2^{31}, 2^{31}-1, -1, 0, 1, 2\}$).

\subsubsection*{7. Common mistakes}

\begin{warningbox}[Mistake 1 --- \texttt{int} slots with an \texttt{INT\_MIN} sentinel]
\begin{lstlisting}[style=snippet]
int a = INT_MIN, b = INT_MIN, c = INT_MIN;          // WRONG: [1,2,INT_MIN] returns 2
long long a = LLONG_MIN, b = LLONG_MIN, c = LLONG_MIN;  // RIGHT
\end{lstlisting}
\end{warningbox}

\begin{warningbox}[Mistake 2 --- no duplicate check]
Without the \texttt{continue}, Example 3 ends with $(a,b,c) = (3, 2, 2)$ and returns $2$; the test $[1,2,2,5,3,5]$ returns $3$ instead of $2$. The check must come \emph{before} the comparisons.
\end{warningbox}

\begin{warningbox}[Mistake 3 --- \texttt{>=} instead of \texttt{>}]
\texttt{if (x >= a)} on a duplicate of $a$ shifts the same value into $b$. With the duplicate check in place the comparison operators are irrelevant for equal values, but writing \texttt{>} keeps the code correct even if someone deletes the check.
\end{warningbox}

\begin{warningbox}[Mistake 4 --- returning the wrong fallback]
With fewer than three distinct values the answer is the \emph{maximum}, not $-1$ and not the second maximum.
\end{warningbox}

\subsubsection*{8. Takeaways}

\begin{itemize}[leftmargin=*]
    \item ``The $k$-th largest'' for a small fixed $k$ is a single-pass aggregate: a sorted list of $k$ slots, updated by insertion.
    \item Decide \emph{distinct values} versus \emph{elements} first; it is a one-line difference in code and a different answer.
    \item When the input may take every value of a type, represent ``empty'' with a wider type, a counter, or a container size --- never with a value of the input type.
    \item A brute force that sorts is often accepted; the single pass is the version that generalises to streams and to $\bigO{1}$ memory.
\end{itemize}

\begin{variantsbox}
\variant{Largest at least twice the others}{\exeasy\ \lc{Easy}\ \varlink{Largest Number At Least Twice of Others}{LeetCode 747}{https://leetcode.com/problems/largest-number-at-least-twice-of-others/}}
  {return the index of the unique largest element if it is at least twice every other element, else $-1$.}
  {the two largest \emph{elements} (duplicates count) and the index of the first; the test is $m_1 \ge 2 m_2$.}
  {one pass with the non-distinct top-2 shift plus an index; $[3,6,1,0] \to 1$, $[1,2,3,4] \to -1$.}
\variant{Product of the two largest}{\exeasy\ \lc{Easy}\ \varlink{Maximum Product of Two Elements in an Array}{LeetCode 1464}{https://leetcode.com/problems/maximum-product-of-two-elements-in-an-array/}}
  {maximise $(a_i - 1)(a_j - 1)$ over $i \ne j$, all $a_i \ge 1$.}
  {non-negative factors: the two largest elements, duplicates allowed.}
  {top-2 shift, answer $(m_1-1)(m_2-1)$; $[3,4,5,2] \to 12$, $[1,5,4,5] \to 16$, $[3,7] \to 12$.}
\variant{Two largest and two smallest}{\exeasy\ \lc{Easy}\ \varlink{Maximum Product Difference Between Two Pairs}{LeetCode 1913}{https://leetcode.com/problems/maximum-product-difference-between-two-pairs/}}
  {maximise $a_w a_x - a_y a_z$ over four distinct indices, all values positive.}
  {the two largest and the two smallest elements: four slots kept in one pass.}
  {answer $m_1 m_2 - s_1 s_2$; $[5,6,2,7,4] \to 34$, $[4,2,5,9,7,4,8] \to 64$.}
\variant{Neither minimum nor maximum}{\exeasy\ \lc{Easy}\ \varlink{Neither Minimum nor Maximum}{LeetCode 2733}{https://leetcode.com/problems/neither-minimum-nor-maximum/}}
  {distinct values; return any element that is neither the minimum nor the maximum, or $-1$.}
  {any three elements suffice: the middle one of them is neither extreme of the whole array.}
  {$\bigO{1}$: if $n < 3$ return $-1$, else the median of the first three; $[3,2,1,4] \to 2$ (any non-extreme value is accepted), $[1,2] \to -1$, $[2,1,3] \to 2$.}
\variant{$k$-th largest distinct value}{\varnone}
  {return the $k$-th largest distinct value, $k$ given (fallback: the maximum).}
  {the top-$k$ distinct list of the theory; $k = 3$ is this problem.}
  {\texttt{TopDistinct}$(k)$ or a set of size $\le k$: $\bigO{nk}$ or $\bigO{n \log k}$.}
\end{variantsbox}
